
\section{Nineteenth Sunday after Trinity: Eighteen Years of Sufferings Ended on a Sabbath}

\begin{quote}

Luke 13:10–17. And he was teaching in one of the synagogues on the Sabbath. And behold, there was a woman who had had a spirit of infirmity eighteen years, and she was bowed together and could not at all straighten herself up. But when Jesus saw her, he called to her and said to her: Woman, you are loosed from your infirmity. And he laid his hands on her, and immediately she straightened herself up and praised God. Then the ruler of the synagogue answered, being angry that Jesus had healed on the Sabbath, and said to the people: There are six days on which one ought to work; come therefore on those days and let yourselves be healed, and not on the Sabbath day! Then the Lord answered him and said: You hypocrite! Does not every one of you on the Sabbath loose his ox or his ass from the stall and lead it away to water it? But this woman, who is a daughter of Abraham, whom Satan had bound, think of it, for eighteen years, ought she not to be loosed from this bond on the Sabbath day? And when he said this, all who were against him were put to shame, and all the people rejoiced over all the glorious works that were done by him.

\end{quote}

\bigskip

When a poor woman has suffered for eighteen years, there is no one who wonders that she would gladly be delivered from her sufferings. It is so heavy a fate to bear that, if anyone could understand even a little of it, compassion would surely fill his soul, and this wish would rise within him: Oh, that help might be found for such misery!

Jesus had power to help, and he helped. The sick woman in our text received unexpected help on that Sabbath when Jesus was teaching in the synagogue. And as her body was healed, so too her soul rose up out of the depth of sorrow and despondency, and God’s praise rose up from her heart and mouth. Happy woman, who on that Sabbath found her way to the synagogue and to Jesus.

But is she the only one who suffers, and the only one who is helped? Is there not a multitude, which no one can number, who suffer and suffer from day to day, from year to year? Are there not millions who bear pain and sorrow through a long life, a life that seems longer still because it is so full of suffering? Yes indeed, there are many, many, innumerable ones who suffer, suffer in the body, suffer in the soul, suffer in heart and mind; sometimes for eighteen years, sometimes for twice eighteen.

And yet there are many among these suffering ones who seek no help in their distress; perhaps there are many among them who do not understand their own misery. For the true suffering of the human heart, and the source of all sufferings, is its sin; and sin—who knows it, who seeks deliverance from it? It is a sight fitted to tear a heart asunder with the deepest sorrow, to behold the incomprehensible indifference of men toward sin and its dreadful consequences. They are sick, sick unto death, sick in soul and heart with a mortal sickness, and yet this desperate madness, that they will have no help, will seek no help! Even if now and then they perceive how the worm gnaws at the root of their life, yet they seek no remedy; they only try to hide their death-wound from themselves and others and to continue their accustomed life, while their heart’s blood runs drop by drop, until there is no hope any more.

Oh, that it were possible to cry it into the ear and heart of the secure sinner: You are sick, you are dying; hasten to seek the Physician, the only Physician who can help you, who can help all.

But then there are also those who feel their misery, and for whom sin has become a grievous and pressing burden. They know not only that they are wretched, but they also know whence their misery comes. But even these often go on for many long periods and do not seek the help, do not seek the Physician. They lack courage to go to him, they lack strength to seek him out; they recoil trembling before the Physician, and they remain lying in the sin that torments them with unspeakable agonies.

Is there anyone who has compassion on all these manifold sufferers? Yes, the same one who had compassion on the woman who for eighteen years had been bound by Satan, Jesus has compassion on all those human souls whom Satan has bound. Whether it be for a long time or a short; whether it be eighteen years or twice as long; whether you understand that you are bound, or imagine that you are free. Jesus has compassion on all these suffering human hearts; he would so gladly make them free. He has the power to do so; for through death he has made powerless the one who had the dominion of death, namely the devil, in order to free those, as many as through fear of death were under bondage all their lifetime. A man bound and ill-used by the devil is an object of the heartfelt compassion of God and of our Savior.

Will you then not be helped, when Jesus both can and will help? Have you no compassion on yourself, so that you will seek out the heavenly Physician? Or do you not know where he is to be found? Try there where that daughter of Abraham found him, in the synagogue, where he was teaching, on the Sabbath, when he spoke the Word of God. There, where the Word of God is spoken, there, where the Gospel is preached, there, where the congregation gathers in prayer, there, where the water of Baptism and the wine of the Supper bestow forgiveness of sins, life, and salvation, there he himself is, the living Son of God, personally present; there he will lay his own helping hands on you, and you shall be healed. The sufferings of a long life shall abruptly end, and it shall become a blessed Sunday for you, if you come bowed down by your sins and sorrows to the Lord’s House, and go forth free, saved, escaped from the devil’s bonds by the forgiveness of sins.

Will you not come? You have suffered long enough, you have been sick so many, many years. All can yet become well again on the Lord’s Day, in the Lord’s House. There the fountain of life is opened, there the remedy is to be found, come, come!

Oh, there are so many who neglect the Lord’s House and the Lord’s Day and the Lord’s Word, though it is their own deliverance and salvation that by this they neglect. Yet there is, God be praised, an incomprehensible riches of the proclamation of God’s Word among us, and yet so many who do not come to the rich spring-well. Congregation of God, awake, and let it be a matter of urgent concern to you that the Lord’s House may be filled, filled with sick, suffering human beings who can find remedy for their souls through the Lord’s Word.

And you, servants of the Lord in his House, awake, and become zealous to speak the Word of life to human hearts. There are many who suffer, and you do not know it; there are many who hunger, without your being able to see it; there are many who weep in secret, whose tears are hidden from you; there is many a heart that bleeds from a hidden wound.

Oh, let a whole and untrimmed witness concerning sin and grace sound forth, cutting asunder and refreshing, wounding and healing, slaying and making alive, every Sunday from your pulpits, and then the incomprehensibly blessed thing shall come to pass from Sunday to Sunday by the power of the Gospel of God: the bound shall be loosed, captives shall be set free, the sick shall be healed, the suffering shall be refreshed, the dead shall become alive. Let none be found like that heartless ruler of the synagogue, who was angry that Jesus healed on the Sabbath. Alas, alas, he is not the only one of his kind in the history of the congregation of God. Every heartless and spiritless priest has the same mind, even if he does not have the same words in his mouth.

Think on the poor souls in the devil’s bonds, and may the Lord himself grant to find releasing words from the Lord’s Law and the Lord’s Gospel, to break the bonds and set the bound at liberty!


